%% title_page_emse.tex -- separate metadata file for Editorial Manager
\documentclass[11pt]{article}
\usepackage[utf8]{inputenc}
\usepackage[T1]{fontenc}
\usepackage[margin=1in]{geometry}
\usepackage{lmodern}
\usepackage{url}
\usepackage{hyperref}
\hypersetup{colorlinks=true,urlcolor=blue}

\begin{document}

\begin{center}
{\Large\bfseries Recovery-Channel Prompt Injection in Self-Healing LLM
Agent Frameworks: A Multi-Case Empirical Study\par}
\vspace{1.2cm}
{\large Mingyuan Xie and Zhengxun Wu\par}
\vspace{0.4cm}
School of Information Management\\
Qingdao University of Technology\\
Qingdao, Shandong, China
\end{center}

\vspace{0.8cm}
\noindent\textbf{Corresponding authors and contact information}

\begin{itemize}
  \item Mingyuan Xie (corresponding author):
  \href{mailto:jhj208436@gmail.com}{jhj208436@gmail.com}
  \item Zhengxun Wu (corresponding author):
  \href{mailto:2276815088@qq.com}{2276815088@qq.com}
\end{itemize}

\noindent\textbf{ORCID.} No ORCID identifiers are provided at initial submission.

\noindent\textbf{Funding.} This study received no targeted external funding.

\noindent\textbf{Competing interests.} The authors declare that they have no
competing interests.

\noindent\textbf{Data availability.} The replication package is available at
\url{https://github.com/dadawer-web/ouisani/releases/tag/emse-recovery-channel-v1}.
It contains the derived trial records, source-audit outputs, provenance and
power-analysis artifacts, and scripts used to generate the reported tables.

\noindent\textbf{Acknowledgements.} None.

\end{document}
